\documentclass[sigconf,nonacm]{acmart}
\usepackage{tabularx}
\sloppy

\begin{document}

\title{External Effect Anomalies in AI Agent Workflows}

\author{Artem Trofimov}
\affiliation{%
  \institution{Independent Researcher}
  \city{Berlin}
  \country{Germany}}
\email{artemtrofimov@acm.org}

\author{Boris Novikov}
\affiliation{%
  \institution{Independent Researcher}
  \city{Helsinki}
  \country{Finland}}
\email{borisnov@acm.org}

\begin{abstract}
AI agents act on the world through tools that send messages, charge cards,
and place orders. Such effects generally cannot be rolled back, only
compensated, and agents release them while retrying failed calls,
speculating, and sharing resources with other executions. The mechanisms for
this --- replay, compensation, staged release, coordinated access --- come
from long-running and multilevel transaction processing; what is new is the
boundary, since an agent typically calls third-party tools that declare
little about themselves.

We give an effect-history model that separates what externalized in the world
from what the runtime observed, and a catalog of eight external-effect
anomalies grouped by the three ways transactional control is lost:
uncertainty, workflow structure, and interaction. Read in the other
direction, the catalog yields the boundary capabilities a runtime must
establish to exclude each anomaly, and four limits beyond which no runtime can
provide a general guarantee from black-box tool invocation alone. Current runtimes supply these
capabilities through adapter layers they control, but such declarations stay
runtime-local rather than reusable across a shared interface. A census of
90{,}818 tools exposed by registered Model Context Protocol servers finds
annotation fields carried at scale but expressing little: three signatures cover 77\% of tools,
and only 13.2\% are in a configuration where the destructiveness hint
applies at all.
\end{abstract}

\maketitle

\section{Introduction}

AI agents increasingly take actions that change the outside world: sending
emails, charging cards, placing orders, calling external services. Such
effects often cannot be undone. Take an agent booking a trip on a 1500 EUR
budget. To save time it books the flight (900) and the hotel (700) in
parallel. Both calls succeed, but together they break the budget. The hotel
can be canceled; the flight was non-refundable. Every tool call worked, yet
the workflow left a state in the world that cannot be fully repaired.

Agents make this common because they run steps in parallel (supervisor-worker
decomposition), plan speculatively (tree search, best-of-$k$), execute likely
actions early to hide latency, and recover by retrying after failures. Each
is useful, but each can also produce a new kind of external-effect failure.
We argue these failures form one catalog, with three sources.
\emph{Uncertainty} failures: when the system cannot confirm the outcome of
an effect --- because of a crash, a timeout, or a lost acknowledgment --- and
must act anyway, the same effect may happen
twice, fail to happen while the workflow proceeds as if it had, or be
compensated blindly. \emph{Workflow} failures: an effect may outlive an
aborted workflow, leak before the workflow decides, or be depended on by a
committed branch though it came from an abandoned one. \emph{Interaction}
failures: an effect meets something outside its own execution --- a
concurrent execution writing the same resource, or an outside actor that
observes and reacts to it before it can be undone. Agents show all three,
because an agent is a long-running workflow that acts on tools whose
outcomes it cannot always confirm and shares an open world with others.

We define correctness through a catalog of anomalies, each a forbidden
pattern in the effect history. The intuition: a safe execution is one whose
surviving external effects can still be explained as the intended committed
workflow --- in the spirit of opacity~\cite{guerraoui2008opacity}, but for
effects that cannot be rolled back. Three features shape the setting: there
is no rollback, only optional and fallible compensation; effects may be
released before the workflow decides; and the surrounding world is open and
may react.

The mechanisms agent runtimes use --- replay, retry, compensation, staged
release, coordinated access --- are not new: they come from long-running and
multilevel transaction processing, including sagas, semantic recovery, and
escrow-style concurrency
control~\cite{garcia1987sagas,korth1990formal,elmagarmid1992transaction,weikum2001transactional,oneil1986escrow},
in the anomaly-based style of classical isolation
levels~\cite{berenson1995critique,adya1999}. The saga lineage isolates the
pivot after which effects are no longer compensable; our catalog classifies
the ways an effect can be incorrect at that point, extended to concurrent,
speculative, and nondeterministic workflows. What differs is the boundary:
those models reason over operations whose semantic properties are known to,
assumed by, or supplied to the transaction manager, whereas an agent
typically calls third-party tools whose properties the interface leaves
largely undeclared. Current runtimes therefore supply the properties they
need themselves --- adapter annotations, registered inverses, per-effect
tiers --- out of band and for themselves alone. Better replay and recovery do not close
the gap on their own: the missing piece is a contract layer at the tool
interface.

This paper makes three contributions. First, an effect-history model for
agent workflows whose operations are calls to externally supplied tools
(Section~\ref{sec:model}). Second, a catalog of eight external-effect
anomalies, organized by the three ways transactional control is lost ---
uncertainty, workflow structure, and interaction --- with guarantee profiles
naming the exclusion sets and a reading of current agent runtimes as partial
coverage of it (Section~\ref{sec:catalog}). Third, for each anomaly, the
tool contracts that make excluding it achievable, and four guarantee
boundaries identifying where black-box invocation alone is insufficient
under our model (Section~\ref{sec:contracts}), grounded by a registry-wide census
of MCP tool annotations (Section~\ref{sec:grounding}).

\section{Effect Histories and the Tool Boundary}
\label{sec:model}

Agent systems form a hierarchy of levels, as in multilevel transaction
management, where the guarantees of one level derive from the properties of
the operations below~\cite{weikum1991principles}. We work with a slice of
two adjacent levels $L_i$ and $L_{i+1}$, written L0 and L1. L0 holds the
operations the agent issues as atomic calls --- typically tools behind
external APIs, \texttt{charge()} or \texttt{book\_flight()}; L1 holds the
workflows composed from them, \texttt{BookTrip} or \texttt{HireCandidate}.
A workflow can itself serve as an operation one level up. The history is
defined at L0 over external events and the runtime's observations of them,
of which the runtime generally sees only a projection; effect safety is
defined at L1, against the workflow's decisions.

The slice is relative downwards too: \texttt{book\_flight} is atomic only
from the caller's side, while for its provider it is a workflow over
operations the caller does not see, of unknown depth. When that hidden
structure shows through --- a timeout inside the provider's workflow --- the
caller is left with an outcome it cannot confirm.

The boundary also varies in who owns it. When agent and tools are built
together, the properties of L0 operations are declared by construction, and
a runtime that owns its adapter layer is in a similar position: it can
attach the properties it needs tool by tool. Neither option transfers
cleanly across a shared interface, where a client sees only what the
interface carries and local declarations are not reusable by other
runtimes. The Model Context Protocol
(MCP) is the widely used interface of this kind and is publicly
inspectable, which is why we return to it in Section~\ref{sec:grounding}.
However the properties are obtained, they constrain which effect-safety
guarantees are achievable at L1.

We focus on five properties of L0 operations: idempotence, invertibility,
externalization, determinism, and commutativity. They are independent ---
\texttt{increment} commutes but is not idempotent, \texttt{delete} is
idempotent but does not commute --- and so form a property space rather than
a single reversible/irreversible scale. Commutativity is moreover
\emph{relational}, a property of operation pairs, and often conditional on
state (escrow-style~\cite{oneil1986escrow}); determinism, not usually
treated as an operation contract, matters because black-box LLM-backed tools
can be nondeterministic.

An execution is a history over two layers: what happened outside, and what
the runtime learned about it (Table~\ref{tab:model}). A retry is a second
attempt of the same logical operation; a compensating action externalizes in
its own right; an exogenous effect is one no attempt precedes and no
compensation reaches; and \emph{unknown} means no acknowledgment arrived, so
whether the effect happened is undetermined for the runtime though
determined in the world. We write $externalize(q, e \text{ on } r)$ when the
resource matters. Where the individual attempt is not at issue we
write $effect(\ell)$ and $cmp(\ell)$.

\begin{table}[htbp]
\centering
\renewcommand{\arraystretch}{1.15}
\footnotesize
\caption{Effect-history vocabulary. The upper group records execution and
external events; the lower, runtime observations and derived predicates.}
\label{tab:model}
\begin{tabularx}{\linewidth}{@{}l X@{}}
$attempt(a, \ell)$ & attempt $a$ of logical operation $\ell$ \\
$externalize(q, e)$ & runtime action $q$ produced external effect $e$ \\
$externalize_{\mathit{ext}}(p, x)$ & external party $p$ produced $x$ \\
$cmp(c, a)$ & compensating action $c$ for attempt $a$ \\
$neutralizes(c, e)$ & $c$ neutralized effect $e$ \\
$dep(e_2 \leftarrow e_1)$ & $e_2$ produced having observed $e_1$ \\
$commit(w)$, $abort(w)$ & workflow $w$ is resolved \\[2pt]
$observe(a, s)$ & outcome for $a$: confirmed, failed, unknown \\
$Req(w)$ & operations the commit of $w$ asserts \\
$resolved(w)$ & $w$ has committed or aborted \\
$survives(e)$ & $e$ externalized, never neutralized \\
\end{tabularx}
\end{table}

Keeping the two layers apart matters below: an effect either occurred or did
not, while \emph{unknown} is a state of the runtime, and acting under it is
what several anomalies turn on. Unlike read-write
histories~\cite{adya1999}, effect histories describe what a workflow
externalized, and the anomalies live in the window between externalization
and $resolved$ --- a window those models do not have.

\section{Effect Anomalies}
\label{sec:catalog}

Each anomaly exposes something missing at the agent--tool boundary:
evidence (A1--A3), lifecycle control (A4--A6), or coordination (A7--A8).
The uncertainty anomalies arise when the runtime cannot confirm the outcome
of an effect --- after a crash, a timeout, or a lost acknowledgment --- and
must act anyway. The workflow anomalies arise from the structure of a single
execution: survival under abort, timing relative to its resolution, and
dependency. The interaction anomalies arise when an effect meets something
outside its own execution: another execution writing the same resource (A7),
or an external actor that observes and reacts to it (A8), which need not be
human --- it may itself be an agent workflow at a level our boundary does not
see.

Table~\ref{tab:catalog} states the eight patterns together with the boundary
capability each one requires. The two readings are dual: a forbidden pattern
says what may go wrong, and the capability says what a runtime must be able
to establish from the boundary to exclude it in general. Boundary
capabilities do not themselves exclude anomalies; protocols consume them,
through safe retry, commit gating, compensation, dependency tracking, and
mediated ordering. Realizations are illustrative: equivalent sources of
evidence, control, or coordination satisfy the same capability.

\begin{table*}[htbp]
\centering
\renewcommand{\arraystretch}{1.3}
\small
\caption{External-effect anomalies and the boundary capabilities required
to exclude them in general. Representative realizations are illustrative
rather than prescriptive.}
\label{tab:catalog}
\begin{tabularx}{\textwidth}{@{}l X X X@{}}
\textbf{} & \textbf{Forbidden pattern} & \textbf{Required boundary capability} & \textbf{Representative realization} \\
A1 Duplicated & one logical operation externalizes twice & authoritative convergence on one outcome & logical-operation ID; idempotent re-issue returning the original outcome \\
A2 Missing & commit without a required effect & authoritative outcome; atomic batch & status endpoint; prepare/commit \\
A3 Orphaned & compensation issued under an unknown outcome & outcome resolution before compensating & outcome query; outcome-conditioned compensation \\
A4 Residue & aborted workflow leaves a surviving effect & residue prevention or safe neutralization & compensation declaration; staging \\
A5 Premature & possibly non-surviving effect externalizes before resolution & pre-externalization observation or control & \texttt{quote}/\texttt{dry\_run}; \texttt{hold} with expiry \\
A6 Contaminated & committed effect depends on a non-surviving effect & stable identity for dependencies & stable effect and resource identity \\
A7 Conflicting & unordered non-commuting effects from independent executions & shared-resource coordination & resource scope; commutativity or inverse; mediator \\
A8 Phantom & exogenous consequence of a compensated effect survives & control of external observability & visibility control; mediated observation \\
\end{tabularx}
\end{table*}

Four of the patterns need more than a line. \emph{A3} is unusual in that the
fault is the action itself, not the state it leaves. The runtime compensates
without knowing whether the original attempt externalized:
\begin{verbatim}
attempt(a1, pay_invoice)
observe(a1, unknown)     # outcome not confirmed
cmp(c1, a1)              # compensate blindly
externalize(c1, refund)  # refund takes effect
=> if a1 never externalized, refund is spurious
\end{verbatim}
A1, A2, and A3 are one family: all three stem from an outcome the runtime
cannot confirm and differ only in the reaction --- re-issue, commit, or
compensate. A1 and A2 are visible in the final history; A3 is wrong at the
moment it is done, whatever the outcome turns out to be. The family also
marks where this setting departs from classical transaction processing.
There, an in-doubt outcome is temporary: participants run a recovery
protocol and resolve it. A third-party tool runs no such protocol, so an
unknown outcome stays unknown until the tool itself offers a way out --- a
durable acknowledgment, a status endpoint, or idempotent re-issue. When
agent and tools are developed together the classical picture applies; the
family concerns the third-party boundary.

\emph{A6} is the effect-world analogue of a dirty read. A committed branch
acts on a temporary reservation made by a branch that is later canceled:
\begin{verbatim}
effect(reserve_room R) [b1]
dep(send_invites <- reserve_room R) [b2]
effect(send_invites) [b2]  # invites name R
commit(b2)                 # b2 wins
abort(b1); cmp(reserve_room R)
=> b2 announced R, which no longer survives
\end{verbatim}
The contamination is causal, not about reversibility: the harm --- invitations
already sent naming a room that was released --- would arise even if those
invitations were themselves reversible.

\emph{A7} covers two coding workers pushing to the same branch from the same
base, neither having observed the other. Commutativity \emph{excludes} the
anomaly; invertibility does not exclude it but makes it \emph{repairable},
letting the loser be rolled back and re-applied toward a serializable state;
mediation prevents the race by ordering access. The conflict is beyond
runtime repair when the effects are both non-commuting and irreversible and
no mediator prevented their unordered externalization.

\emph{A8} differs from all of these. An external party acts on an effect
that is later withdrawn:
\par\noindent\begin{minipage}{\linewidth}
\begin{verbatim}
externalize(a1, e)       # offer sent
externalize_ext(S, x)    # supplier acts on it
dep(x <- e)              # x saw e
cmp(c1, a1)              # offer withdrawn
neutralizes(c1, e)
survives(x)              # supplier action stands
=> offer retracted; the reaction is not
\end{verbatim}
\end{minipage}
The dependent effect $x$ is exogenous: no attempt of the runtime produced it
and no compensation reaches it. This separates A8 from A6, where the
surviving effect is managed. A1--A7 concern effects within the managed or
mediated boundary and can be detected and sometimes repaired after the fact;
an exogenous consequence lies outside the projection the runtime observes
and can compensate. The reacting party may also sit outside any contractual
boundary, so no ordering or commutativity can be declared for it, and the
enforceable primitive is preventive: control whether and when the effect
becomes externally observable.

A5 and A7 are profile-relative preventive patterns: a particular execution
may happen to end safely --- the resource may serialize the calls on its own,
and an early effect may turn out to survive --- but the stronger profile
forbids relying on an outcome the boundary does not establish.

\begin{table*}[htbp]
\centering
\renewcommand{\arraystretch}{1.3}
\small
\caption{Coverage of the anomaly catalog by current agent runtimes, under
the assumptions each system states.}
\label{tab:systems}
\begin{tabularx}{\textwidth}{@{}l X X@{}}
\textbf{System} & \textbf{Targeted coverage} & \textbf{Notes} \\
\hline
ACRFence~\cite{zheng2026acrfence} & unknown-safe partial (A1) & proposed; attack validated, not implemented \\
RAC~\cite{perera2026rac} & compensation-safe partial (A4) & no gating; no idempotency, so A1--A3 not addressed \\
Atomix~\cite{mohammadi2026atomix} & A1/A3 partial, compensation-safe (A4), speculation-safe (A5, A6), externally-mediated partial (A7) & idempotent-known-outcome path; no crash-safe exactly-once; A2 gap on multi-effect release \\
Cordon~\cite{chen2026cordon} & A1/A3 partial, compensation-safe (A4), speculation-safe partial (A5) & idempotent tools; staged effects; task-scoped \\
CoAgent~\cite{lyu2026coagent} & externally-mediated partial (A7) & multi-agent serializability via reorder + registered inverse; irreversibles gated \\
Shepherd~\cite{yu2026shepherd} & per-effect reversibility tiers; A7 observed in supervisor use case & meta-agent substrate, not a runtime guarantee; irreversibles recorded for audit \\
\end{tabularx}
\renewcommand{\arraystretch}{1}
\end{table*}

\subsection{Safety Guarantee Profiles}
\label{sec:levels}

The catalog induces a vocabulary of \emph{effect-safety guarantees}, each
naming a set of anomalies it excludes: \textbf{unknown-safe} (A1--A3),
\textbf{compensation-safe} (A4, assuming compensations are correct and
succeed), \textbf{speculation-safe} (A5--A6), and
\textbf{externally-mediated} (A7--A8, where the boundary admits mediation).

Interactive tasks explain why compensation-safe is sometimes the best
achievable profile: an agent negotiating on a user's behalf cannot learn the
counterparty's reaction without sending an offer, so the action must precede
the observation and full gating is unavailable. Where the inputs needed to
resolve the workflow can be observed without effect --- through a
\texttt{quote} --- the stronger profile becomes attainable.

\subsection{Current Runtimes Against the Catalog}
\label{sec:systems}

Existing systems target different fragments of the catalog under
runtime-specific assumptions. Table~\ref{tab:systems} summarizes their
stated coverage rather than proven guarantees; none realizes a clean
cumulative hierarchy.


Atomix and Cordon stage effects, with Atomix additionally tracking
speculative dependencies; their A1/A3 coverage is conditional on assumed
tool semantics rather than on a general outcome-resolution protocol. CoAgent
reorders shared-resource effects using registered inverses and gates
irreversible calls when no inverse exists --- an instance
of the achievability condition for A7 rather than an unconditional
guarantee. Shepherd exposes reversibility tiers as a substrate rather than a
guarantee. Across these systems, coverage rests on runtime-owned
declarations --- adapter-declared idempotency keys, reversibility
annotations, registered inverses, per-effect tiers --- rather than a
reusable shared contract; none controls exogenous reactions (A8), and no
system covers the full catalog. Adjacent work governs the \emph{read} side
of shared agent state~\cite{khan2026sbus} and intent revision under
conflicting irreversible actions~\cite{zhai2026revisable}, applies
Saga-style validation to multi-agent planning~\cite{chang2025sagallm}, and
argues for isolation-level reasoning in workflow
systems~\cite{stonebraker2025consistency}.

\subsection{Scope and Coverage}
\label{sec:coverage}

The catalog is coverage-oriented rather than complete by theorem: absolute
completeness is unavailable in an open world, where an external actor can
react in unbounded ways. Each family has a structural reason for its
membership. The uncertainty family is fixed by the reactions available when
an outcome is not confirmed: re-issue, commit as though it had
externalized, or compensate. The workflow family is fixed by how an
externalized effect relates to its resolution point --- survival, timing,
dependency. The interaction family separates unmediated concurrent effects
on a shared resource from an exogenous consequence of a later-compensated
effect. We conjecture that, relative to the vocabulary of
Section~\ref{sec:model}, every loss of transactional control the model can
represent falls into one of the three families; formal coverage,
minimality, and independence are future work.

We place outside the model, so the catalog is not mistaken for covering
them: semantic planning errors (the agent booked the wrong city); read-side
anomalies, the subject of read-side isolation work; security and policy
violations; tool-contract misclassification (a tool declared reversible that
is not); intra-execution effect order, except insofar as it participates in an A7
conflict; and liveness failures.

\section[Transactional Tool Contracts and Guarantee Boundaries]{Transactional Tool Contracts and\\ Guarantee Boundaries}
\label{sec:contracts}
\label{sec:unknown}

Existing runtimes obtain tool semantics through adapter layers they control.
Classical transaction models assume properties such as idempotence,
invertibility, and commutativity are known; at a third-party boundary they
are often absent or merely inferred. A reusable contract layer must
therefore distinguish established, assumed, and unknown properties, since
acting on an assumption --- compensating an operation with no true inverse
--- is a failure mode of its own. Such declarations do not replace runtime
protocols: they are their precondition.

A runtime consuming tools through a shared interface reads only what that
interface carries, and MCP is not entirely silent: since its 2025-03-26
revision, tool annotations let a server flag a tool as read-only,
destructive, idempotent, or open-world.\footnote{Model Context Protocol
specification, revision 2025-03-26.} These are advisory boolean hints ---
unenforced, optional, and coarser than the capabilities of
Table~\ref{tab:catalog}: \texttt{idempotentHint} marks a property but
supplies no idempotency key, and no hint covers status resolution,
compensation, staging, commutativity, or visibility. How servers populate
these fields is measured in Section~\ref{sec:grounding}.

\subsection{Recurring Contract Families}

The capabilities in Table~\ref{tab:catalog} fall into three contract
families. A capability is what the boundary lets a runtime establish; a
contract is the reusable declaration that exposes it, possibly composed of
several operations and their documentation.

Uncertainty requires authoritative convergence on one logical outcome: a
status query, a durable acknowledgment, or idempotent re-issue returning the
original result. Deduplication alone excludes A1 but does not resolve the
outcome that A2 and A3 turn on. The adapter-declared idempotency key is the
closest primitive in use among the systems of Section~\ref{sec:systems};
atomic multi-tool release additionally requires prepare/commit
participation, which none of them supplies.

Lifecycle control (A4--A6) needs the workflow's invariants to be checkable
before irreversible effects occur, compensation semantics for what a
compensator requires and what counts as its success, and stable effect and
resource identity afterwards. A read-only \texttt{quote} or an expiring
\texttt{hold} provides pre-externalization control; the unit dependency
tracking protects is a \emph{commit sphere}, the minimal set of effects that
must become final together.

Coordination (A7--A8) requires contracts at a different granularity.
Idempotency, status, quote, and compensation are properties of a single
operation, whereas commutativity holds of a \emph{pair} of operations on a
resource and cannot be set by rules given to one agent: it must be declared
at the resource, together with inverses and, where neither holds, a mediator
that serializes access. Visibility is likewise a property of the boundary:
who may observe an effect before it is final.

A boundary that exposes these capabilities makes the corresponding exclusions
implementable by any runtime, given a suitable protocol and declarations
that hold, in place of the out-of-band declarations of
Section~\ref{sec:systems}.

\subsection{Guarantee Boundaries Above Black-Box Tools}
\label{sec:boundaries}

Four boundaries mark where black-box invocation alone is insufficient:
stronger guarantees require authoritative tool participation, shared
mediation, or visibility control. We state them as consequences of the
definitions, instantiated by current systems; formal proofs are future
work.

First, without a primitive that lets retries converge on the authoritative
outcome --- idempotent re-issue returning the original outcome, a status
endpoint, or a durable acknowledgment --- no general unknown-safe execution
is available: over an unreliable channel the runtime cannot learn in bounded
time whether an effect took place, the classical exactly-once
barrier~\cite{bernstein1987concurrency} at the agent--tool boundary, and can
only choose which reaction to risk, trading A1 against A2 against A3.

Second, non-commuting irreversible effects without a mediator admit no
general conflict repair, since repair means restoring a serial order by
undoing and re-applying effects, which irreversibility rules out. What
remains is prevention: mediation before unordered externalization, or gating
conflicting calls until earlier ones commit, as CoAgent does.

Third, an open-world reaction ends the reach of compensation: once an outside
actor has observed a released effect and acted, retracting the effect does
not retract the reaction (A8). The available contract is preventive ---
visibility control before observation.

Fourth, several irreversible effects on different tools cannot be released
atomically above the tool layer; a crash between releases leaves some
effects placed and others missing (A2). Making the batch atomic is the
classical atomic-commitment problem~\cite{bernstein1987concurrency} and
needs tool-side prepare/commit participation, a limit Atomix states for its
own release step.

\section{What the Boundary Declares Today}
\label{sec:grounding}

We census what a runtime can establish from the standard MCP annotation
vocabulary. The measurement covers declarations at the agent-facing
boundary, not service behavior: a value states what a server asserts, and we
do not verify it. A capability we do not find is therefore \emph{not
established from the annotations}, and may still follow from a tool's input
schema, its documentation, or the upstream API it wraps.

We took a full snapshot of the official MCP registry on 2026-07-28
(59{,}625 entries; 18{,}688 distinct servers, keeping the latest version of
each) and queried every anonymously reachable remote server with
\texttt{tools/list} (15\,s timeout, no retries), against specification
revision 2025-03-26. No tool was
called. Of 7{,}987 such targets, 4{,}584 returned at least one tool,
yielding 90{,}818 tools (median 11 per server). The remainder failed to
connect or timed out; servers requiring authentication and stdio-only
packages are out of scope, so the sample covers anonymously reachable
remote integrations. For each tool we recorded the four annotations the
specification defines --- \texttt{readOnlyHint}, \texttt{destructiveHint},
\texttt{idempotentHint}, \texttt{openWorldHint} --- keeping an explicit
\texttt{false} distinct from an omitted field.

\paragraph{Annotations are widely transmitted.}
76.3\% of tools carry at least one annotation and 65.3\% carry all four.
Transmission is not the same as
semantic information, however: the four booleans are serialized whether or
not their values say anything about the particular tool. The specification
renders \texttt{destructiveHint} vacuous when \texttt{readOnlyHint} is
true, and that pairing alone covers 56\% of all tools.

\paragraph{The values that are sent are concentrated.}
Reading each tool as a signature over the four fields --- each true, false,
or omitted --- 66 of the 81 possible signatures appear, but their
frequencies are concentrated: one signature (read-only, non-destructive,
idempotent, open-world) covers 42.7\% of tools, and three cover 76.7\%. 77.3\% of annotating multi-tool servers use more than one signature, so the
concentration is not one default copied across every tool;
nevertheless, the effective annotation vocabulary a runtime encounters is
dominated by a few recurring patterns. The limitation is not presence on the wire but the resolution of
what the fields can say; we do not attribute the values to author intent,
since many may originate in SDK defaults or server templates.

\paragraph{What the hints do not resolve.}
The specification makes \texttt{destructiveHint} meaningful only when
\texttt{readOnlyHint} is false. The field is present on 69.3\% of tools,
but only 13.2\% are in a position to say anything about destructiveness at
all, and 3.0\% assert a destructive operation.

\medskip

\noindent
The four hints express only coarse properties of a call: whether it may
modify its environment, whether such modification may be destructive,
whether repetition adds further effects, and whether it interacts with
external entities.
Table~\ref{tab:gap} reads the boundary capabilities of
Section~\ref{sec:contracts} against that vocabulary.

\begin{table}[htbp]
\centering
\renewcommand{\arraystretch}{1.2}
\footnotesize
\caption{Expressibility of the boundary capabilities of
Section~\ref{sec:contracts} in current MCP tool annotations.}
\label{tab:gap}
\begin{tabularx}{\linewidth}{@{}X l@{}}
\textbf{Boundary capability} & \textbf{In annotations} \\
A1 authoritative convergence & Limited: idempotence only \\
A2 authoritative outcome; atomic batch & No \\
A3 outcome resolution before compensating & No \\
A4 residue prevention or neutralization & No \\
A5 pre-externalization observation/control & No \\
A6 stable identity for dependencies & No \\
A7 shared-resource coordination & No: externality only \\
A8 control of external observability & No \\
\end{tabularx}
\end{table}

\noindent
The gap is not confined to effectful tools --- read-only operations also
have properties the vocabulary cannot state, such as determinism or
freshness --- but these are the ones a transactional protocol consumes.

Current MCP metadata is thus sufficient to distinguish coarse execution
modes --- read against write, declared idempotent versus potentially
non-idempotent --- and
insufficient to express the semantics a runtime needs to reason about
retries, compensation, speculation, or concurrency. The ecosystem already
provides metadata at scale; what is missing is the vocabulary needed for
transactional reasoning.

\section{Conclusion and Future Work}

We identify eight recurring structural ways an agent workflow loses
transactional control over external effects, from three sources.
Each anomaly identifies a boundary capability --- evidence, control, or
coordination --- that a runtime needs in order to exclude it generally, and current runtimes cover parts of the
catalog by supplying those declarations themselves, out of band. Better replay and recovery do not close this gap on their
own: the missing layer is a set of transactional contracts at the tool
interface.

Future work includes recovering how much of each capability is available
from server implementations and upstream
APIs~\cite{toeppe2026mcpdataset}, and where it is lost at the wrapper;
formalizing the guarantee boundaries; synthesizing protocols from declared
contracts; and studying coverage of the catalog and composition across more
than two levels, where a workflow's contract becomes an operation contract
one level up.

\bibliographystyle{ACM-Reference-Format}
\bibliography{../../bibliography/flame-stream}

\end{document}